% Results per research question (report plan). Every number comes from a file
% in results/ or notes/, never typed in from memory.
\section{Results}

\subsection{RQ1: What can pattern matching detect reliably?}
\pending{table: precision, recall, and F1 with intervals per type; confusion matrix; timing. Then the field set.}

\subsection{RQ2: Internal or external threats?}
\textbf{(a) Published data.} The 2026 DBIR attributes 88\% of breaches with a known actor to external actors and 12\% to internal ones \cite{dbir2026}. The internal share was 35\%, 18\%, and 12\% in the 2024, 2025, and 2026 editions \cite{dbir2024,dbir2025,dbir2026}; the DBIR links such swings mainly to changes in its data contributors, and most internal cases are errors rather than deliberate misuse. Internal shares are higher in some sectors, up to 44\% in public administration \cite{dbir2026}. The DBIR describes its data as a sample of convenience rather than a representative one \cite{dbir2026}, so these figures show which threat is more often observed in reported breaches, not its true prevalence.

\textbf{(b) Our audits, as a case study.} \pending{findings per attacker profile (P1 to P4) and severity, from results/findings/.}

\subsection{RQ3: The most vulnerable part of the system}
\pending{findings grouped by component and severity; restore drill times as supporting evidence.}
